\newcommand*{\path}{chapters/chapter2}

\section{Organic Structures}


\subsection{Hydrocarbon frameworks and functional groups}

Organic chemistry is the study of molecules containing carbon. These molecule also contain hydrogen and sometimes oxygen, nitrogen and other elements.

\marginnote{Amino acids contain two functional groups: an amino (\ce{NH2} or \ce{NH}) group and a carboxylic acid (\ce{COOH}) group}[0cm] 
The chemistry of organic molecules depend less on the number and arrangement of hydrogen and carbon atoms than on other types of atoms. The parts of an organic molecule containing one or more other atoms are called \definition{functional groups}.
\begin{itemize}
    \item \key{The functional groups determine the way the molecule works both chemically and biologically.}
    \item \key{The hyrdocarbon framework is made up of chains and rings of carbon atoms, and it acts as a support for the functional groups.}
\end{itemize}


\subsection{Drawing molecules}

Guidelines for drawing molecules:
\begin{itemize}
    \item \key{Draw chains of atoms as zig-zags.}
    \item \key{Miss out the \ce{H}s attached to carbon atoms along with the \ce{C-H} bonds.}
    \item \key{Miss out the capital \ce{C}s representing carbon atoms.}
\end{itemize}
These are guidelines, they are not rules. If a carbon is written as \ce{C}, add all \ce{H} atoms must be written too.

It's possible to draw the same molecule in different ways on different occasions, in order to emphasize different points.


\subsection{Hyrdocarbon frameworks}
 
Carbon is unusual since it forms strong, stable bonds with the majority of elements in the periodic table (including itself).

\subsubsection*{Carbon chains}

The simplest class of hydrocarbon frameworks contains just chains of atoms (see Table \ref{tab:Cchanis}).

\input{\path/tables/carbon_chains.tex}

\subsubsection*{Organic elements}

\marginnote{Amino acid:\\ \includegraphics[width=\marginparwidth]{\path/molecules/amino_acid}}[0cm]
\ce{R} in a structure can mean anything, is a wild card. The reactivity of organic molecules is so dependent on their functional groups that the rest of the molecule can be irrelevant and we can just choose to call it \ce{R}.

 Abbreviations for the names of carbon chains (Table \ref{tab:Cchanis}) looks like an element symbol: these symbols are sometimes called \definition{organic elements}. These symbols (and names) can be used only for terminal chains of atoms.

\subsubsection*{Carbon rings}

\marginnote{Phenyl group:\\ \includegraphics[width=\marginparwidth]{\path/molecules/phenyl_group}}[0cm] 
Rings of atoms are also common in organic structures. When a benzene ring is attached to a molecule by only one of its carbon atoms it is called a \definition{phenyl group} (with organic element symbol \ce{Ph}). Any compound containing a benzene ring or a related ring system is known as \definition{aromatic}. An \emph{aryl group} (with organic chemistry symbol \ce{Ar}) represents any substituted phenyl group, i.e. a phenyl group with any number of the hydrogen atoms replaced by other groups. Like \ce{R}, the wild card alkyl group, \ce{Ar} is a wild card aryl group.

Rings of carbon atoms are given names starting with \definition{cyclo}, following by the name for the carbon chain containing the same number of carbon atoms.

\subsubsection*{Branches}

Hydrocarbon frameworks are often branched and some short branched carbon chains are given names and organic element symbols. The most common is the isopropyl group (\ce{i-Pr}).
\begin{center}
    \includegraphics[width=\marginparwidth]{\path/molecules/isopropyl}
\end{center}

\key{Isomers are molecules with the same kinds and numbers of atoms joined up in different ways.} Isomers need not have the same functional groups.

The isobutyl group (\ce{i-Bu}) is a \ce{CH2} group joined together with an \ce{i-Pr} group. The \textit{sec}-butyl group (\ce{s-Bu}) has a methyl and an ethyl group joined to the same carbon atom. The \textit{tert}-butyl group (\ce{t-Bu}) has three methyl groups joined to the same carbon atom.

\input{\path/figures/butyl_branches.tex}

\key{The prefixes \textit{sec} and \textit{tert} are short for secondary and tertiary, terms referring to the carbon atom that attaches these groups to the rest of the molecule. A primary carbon atom is attached to only one other \ce{C} atom, a secondary to two other \ce{C} atoms, and so on.} This means there are five types of carbon atom.


\subsection{Functional groups}

The understanding of functional groups is the key to the understanding of organic chemistry.

\subsubsection*{Alkanes contain no functional group}

The alkanes are the simplest lass of organic molecules because they contain no functional groups. They are extremely unreactive.

\subsubsection*{Alkenes (or olefins) contain \ce{C=C} double bonds}

\ce{C=C} double bonds impact the reactivity of an organic molecule as any other functional groups, therefore they are classified as functional groups.

\subsubsection*{Alkynes contain \ce{C#C} triple bonds}

\ce{C#C} triple bonds have a special type of reactivity associated with them, so it's useful to call a \ce{C#C} triple bond a functional group. Alkynes are linear so they are drawn with four carbon atoms in a straight line.

\subsubsection*{Alcohols (\ce{R-OH}) contain a hydroxyl (\ce{OH}) group}

Molecules containing hydroxyl groups are often soluble in water, and living organisms often attach sugar (rich in hydroxyl groups) to otherwise insoluble compounds to keep them in solution in the cell.

\subsubsection*{Ethers (\ce{R^1 -O-R^2}) contain an alkoxy group (\ce{-OR})}

\marginnote{THF:\\ \includegraphics[width=\marginparwidth]{\path/molecules/THF}}[0cm] 
The name ethers refers to any compound that has two alkyl groups linked together through an oxygen. Tetrahydrofuran (THF) is a commonly used solvent and is a cyclic ether.

\subsubsection*{Amines (\ce{R-NH2}) contains the amino (\ce{NH2}) group}

The amino group is present in amino acids. Amines often have a powerful fishy smell. Many neurologically active compounds are also amines (like amphetamine).